% arara: xelatex
\documentclass[12pt]{article}

\usepackage{physics}
\usepackage{microtype}
\usepackage{array}
\usepackage{indentfirst}
\usepackage{sectsty}
\allsectionsfont{\centering}
\usepackage{amsmath, amsfonts, amssymb}
\usepackage{comment}
\usepackage[top=2cm, left=1.2cm, right=1.2cm, bottom=2cm]{geometry}
\usepackage{lastpage}
\usepackage{enumitem}
\usepackage{caption}
\usepackage{url}
\usepackage{fancyhdr}
\usepackage{booktabs}
\usepackage{fontspec}
\usepackage{polyglossia}

\setmainlanguage{russian}
\setotherlanguages{english}
\setmainfont{Arial}
\setlength{\headheight}{15pt}

\pagestyle{fancy}
\lhead{Временные ряды}
\chead{}
\rhead{Контрольная работа}
\lfoot{}
\cfoot{DON'T PANIC}
\rfoot{\thepage/\pageref{LastPage}}
\renewcommand{\headrulewidth}{0.4pt}
\renewcommand{\footrulewidth}{0.4pt}

\AddEnumerateCounter{\asbuk}{\russian@alph}{щ}
\setlist[enumerate, 2]{label=\asbuk*),ref=\asbuk*}

\DeclareMathOperator{\Cov}{\mathbb{C}ov}
\DeclareMathOperator{\Corr}{\mathbb{C}orr}
\DeclareMathOperator{\Var}{\mathbb{V}ar}
\let\P\relax
\DeclareMathOperator{\P}{\mathbb{P}}
\DeclareMathOperator{\E}{\mathbb{E}}
\DeclareMathOperator{\plim}{plim}

\newcommand \hb{\hat{\beta}}
\newcommand \hs{\hat{\sigma}}
\newcommand \htheta{\hat{\theta}}
\newcommand \s{\sigma}
\newcommand \hy{\hat{y}}
\newcommand \hY{\hat{Y}}
\newcommand \e{\varepsilon}
\newcommand \he{\hat{\e}}
\newcommand \z{z}
\newcommand \hVar{\widehat{\Var}}
\newcommand \hCorr{\widehat{\Corr}}
\newcommand \hCov{\widehat{\Cov}}
\newcommand \cN{\mathcal{N}}
\newcommand \RR{\mathbb{R}}
\newcommand \NN{\mathbb{N}}
\newcommand{\cF}{\mathcal{F}}

\begin{document}

\begin{enumerate}
    
    \item \textbf{(10 баллов)} Рассмотрим модель $ETS(AAN)$:
    \[
    \begin{cases}
    u_t \sim \cN(0;\sigma^2),\\
    \ell_t = \ell_{t-1} + b_{t-1} + \alpha u_t,\\
    b_t = b_{t-1} + \beta u_t,\\
    y_t = \ell_{t-1} + b_{t-1} + u_t.
    \end{cases}
    \]
    
    \begin{enumerate}
        \item (5 баллов) Выразите $y_t$ через $y_{t-1}$, $y_{t-2}$ и ошибки $u_t$, $u_{t-1}$, $u_{t-2}$ (то есть получите ARIMA-представление модели).
        
        \item (5 баллов) Найдите математическое ожидание и условную дисперсию
        \[
        \E(y_{t+1}\mid \mathcal{F}_t), \quad \Var(y_{t+1}\mid \mathcal{F}_t).
        \]
        
    \end{enumerate}
    
    \item \textbf{(15 баллов)} Доходности актива описываются моделью $GARCH(1,1)$:
    \[
    r_t = \sigma_t z_t, \qquad z_t \sim \cN(0;1),
    \]
    \[
    \sigma_t^2 = 1 + 0.25\,r_{t-1}^2 + 0.6\,\sigma_{t-1}^2.
    \]
    
    \begin{enumerate}
        \item (5 баллов) Существует ли у процесса безусловная дисперсия? Если да, найдите $\Var(r_t)$.
        \item (5 баллов) Пусть $r_T=2$ и $\sigma_T^2=5$. Найдите
        \[
        \E(\sigma_{T+1}^2\mid \mathcal{F}_T)
        \quad \text{и} \quad
        \E(\sigma_{T+2}^2\mid \mathcal{F}_T).
        \]
        \item (5 баллов) Найдите
        \[
        \lim_{h\to\infty}\E(r_{T+h}^2\mid \mathcal{F}_T).
        \]
    \end{enumerate}
    
    \item \textbf{(15 баллов)} Рассмотрим уравнение
    \[
    (1-0.7L)(1-L)y_t = 3 + (1+0.4L)\e_t,
    \qquad \e_t \sim WN(0,1).
    \]
    
    \begin{enumerate}
        \item (5 балла) Классифицируйте этот процесс как $ARIMA(p,d,q)$. Укажите $p,d,q$.
        \item (5 балла) Найдите первые два значения ACF и PACF после взятия d
        \item (5 балла) Найдите предел после взятия d
        \[
        \lim_{h\to\infty}\E(y_{T+h}-y_{T+h-1}\mid \mathcal{F}_T).
        \]
    
    \end{enumerate}
    
    \item \textbf{(15 баллов)} $(y_t)$ описывается уравнением
    \[
    y_t - 0.9 y_{t-1} + 0.2 y_{t-2} =5 + u_t - 0.5 u_{t-1},
    \qquad u_t \sim WN(0,1).
    \]
    
    \begin{enumerate}
        \item (5 баллов) Разложите лаговые полиномы на множители
        \item (5 баллов) Выпишите более простое уравнение с тем же множеством стационарных решений
        \item (5 баллов) Для стационарного решения найдите явный вид $MA(\infty)$-представления.
    \end{enumerate}
    
\item \textbf{(10 баллов)} Для двух функций $\gamma^{(1)}_k$ и $\gamma^{(2)}_k$, заданных на $\mathbb{Z}$, известно, что
\[
\gamma^{(j)}_{-k}=\gamma^{(j)}_k, \qquad j=1,2.
\]

Кроме того, для $k \ge 0$:
\[
\gamma^{(1)}_k =
\begin{cases}
1, & k=0,\\
0.8, & k=1,\\
0.2, & k=2,\\
0, & k\ge 3,
\end{cases}
\]

\[
\gamma^{(2)}_k =
\begin{cases}
1, & k=0,\\
0, & k=1,\\
\frac12, & k=2,\\
0, & k\ge 3.
\end{cases}
\]

\begin{enumerate}
    \item Для каждой из функций $\gamma^{(1)}_k$ и $\gamma^{(2)}_k$ проверьте, может ли она быть ковариационной функцией некоторого стационарного процесса.
    
    \item (5 баллов) Для каждой функции, которая может быть ковариационной функцией стационарного процесса, предложите пример соответствующего процесса.
\end{enumerate}
    
    \item \textbf{(15 баллов)} Пусть $(\e_t)$ --- последовательность независимых случайных величин, принимающих значения $-1$ и $1$ с вероятностями $1/2$ и $1/2$. Рассмотрим процесс
    \[
    x_t = \e_t \e_{t-1}.
    \]
    
    \begin{enumerate}
        \item (5 баллов) Является ли процесс $(x_t)$ слабо стационарным?
        \item (5 баллов) Является ли он белым шумом?
        \item (5 баллов) Являются ли величины $x_t$ независимыми? Ответ обоснуйте.
    \end{enumerate}
    
\end{enumerate}



\end{document}
